\documentclass[11pt]{article}
\usepackage{../../styles/activity}

\usepackage{xr}
\externaldocument{0-MR}

\lhead{}
%\chead{\textbf{\Large{\hspace{0pt}Beginning Activities for Section~3.2}}\\\prevhead}
\bahead{3.2}
\rhead{}
\lfoot{}
\rfoot{}
\cfoot{\hspace{0pt}\scalebox{0.4}{\includegraphics{cc-by-nc-sa.eps}}}

\begin{document}
\input table
\subsection*{Beginning Activity 1 (Using the Contrapositive)}
There are statements in mathematics for which it is difficult (if not impossible) to construct a direct proof.  In these situations, we frequently use some of the logical equivalencies we studied in Section~2.2 to write a statement that is equivalent to the given statement.  One of the ones used most often is the logical equivalence between a conditional statement and its contrapositive.  The idea is that if we prove one of these to be true (or false), the other one is then true (or false).  An example of this is explored in this beginning activity.

\begin{enumerate}
  \item The proposition appears to be true since for each example,   $n$  is odd whenever  $n^2$  is odd.
  \item It is really not possible to construct a now-show table for a direct proof.  Hopefully, you discovered this.  If we assume that  $n^2 $ is even, we can conclude that there exists an integer  $k$  such that  $n^2  = 2k$.  

It would seem that we could use a square root to obtain something about  $n$, but there are two problems with this.  One is that the square root of  $n^2 $ is equal to the absolute value of  
$n$, and the other is that we cannot really do anything algebraically with  $\sqrt {2k} $.  That is, we cannot really simplify the equation  $\sqrt {n^2 }  = \sqrt {2k} $.

Another possibility could be to divide both sides of the equation  $n^2  = 2k$ by $n$.  One problem with this is that we first have to establish that  $n \ne 0$.  In the situation where  $n \ne 0$, we would get  $n = \dfrac{{2k}}{n}$.  The problems here are that the integers are not closed with respect to division and this does not really help us prove that  $n$  is odd.

  \item The contrapositive is:  For each integer $n$, if $n$ is an even integer, then $n^2$ is an even integer.

\item 
\begin{tabular}[t]{| c | l | l |} \hline
\textbf{Step} & \textbf{Know} & \textbf{Reason} \\ \hline
$P$  &  $n$  is an even integer.  & Hypothesis \\ \hline
$P1$  &  $( {\exists q \in \mathbb{Z}} )( {n = 2q} )$ & Defintion of even integer \\ \hline
$P2$  &  $n^2  = ( {2q} )^2 $ &  Square both sides of the equation \\ \hline
$P3$  &  $n^2 = 4 ( q^2 ) = 2 ( 2q^2 ) $  &  Algebra \\ \hline
$P4$  &  $2q^2$ is an integer.  &  Closure Property of the integers \\ \hline
$Q1$  &  There exists an integer $k$ such that $n^2 = 2k$.  &  The integer is $2q^2$. \\ \hline
$Q$   &  $n^2$ is an even integer.  &  Defintion of even integer \\ \hline
\textbf{Step} & \textbf{Show} & \textbf{Reason} \\ \hline
\end{tabular}

\item Since the contrapositive is logically equivalent to the conditional statement, by proving the contrapositive, we have also proven the original proposition.

\end{enumerate}
\hbreak



\newpage
\subsection*{\textbf{Beginning Activity 2 (A Biconditional Statement)}}


\begin{enumerate} %\setcounter{enumi}{1}
  \item 
$$
\BeginTable
\BeginFormat
| c | c | c | c | c | c |
\EndFormat
\_6
| $P$  |  $Q$    \|6  $P \to Q$  |  $Q \to P$  |  $(P \to Q) \wedge (Q \to P)$ | $P \leftrightarrow Q$ | \\+22 \_6
          | T | T \|6 T | T | T | T | \\ 
         | T | F  \|6 F | T | F | F | \\ 
          | F | T \|6 T | F | F | F | \\ 
          | F | F \|6 T | T | T | T |  \\ \_6
\EndTable
$$

\item We can prove the biconditional statement  $P \leftrightarrow Q$  by completing a proof for  $P \to Q$  and completing a proof for  $Q \to P$.

\item In this case, if use  $P \to Q$  to represent, ``If  $n$  is an odd integer, then  
$n^2 $ is an odd integer'', then  $Q \to P$  would represent,  ``If  $n^2$ is an odd integer, then  $n$  is an odd integer.''  The logical equivalency in Part~(1) tells us that if we have proven the conditional statement  $P \to Q$  and have proven its converse  $Q \to P$, then we have proven the biconditional statement  $P \leftrightarrow Q$.  In this case,  the biconditional statement  
$P \leftrightarrow Q$ is, ``The integer  $n$  is an odd  if and only if  $n^2 $ is an odd integer.''

\end{enumerate}
\hbreak


\end{document}

